The cross partition $\mathcal{R} \otimes \mathcal{S}$ of two partitions $\mathcal{R}$ and $\mathcal{S}$ is the family of sets defined by

$$\mathcal{R} \otimes \mathcal{S} = \{R_i \cap S_j \mid R_i \in \mathcal{R}, S_j \in \mathcal{S}, R_i \cap S_j \neq \emptyset\}.$$

For example,

Since the intersection of measurable sets is measurable and since the sets $\{R_i \cap S_j\}$ are disjoint and exhaustive, a cross partition is a partition.

In Example 2 of Section 2-7 we noted that every denumerable-valued random variable determines a natural partition of $\Omega$. We call the partition associated with the denumerable-valued random variable $\mathbf{g}$, $\mathcal{R}^g$. In terms of the natural partition induced by $\mathbf{g}$, we define the conditional mean of $\mathbf{f}$ given $\mathbf{g}$, by
